\begin{table}[H]
\centering
\caption{Standardized Effect Sizes for Main Outcomes}
\label{tab:sde}
{\small
\begin{tabular}{lcccccl}
\toprule
Outcome & $\hat{\beta}$ & SD($X$) & SD($Y$) & SDE & SE(SDE) & Classification \\
\midrule
House Prices (TCJA) & $-0.0322$ & 0.810 & 0.666 & $-0.0391$ & 0.0098 & Small negative \\
\bottomrule
\end{tabular}}
\par\vspace{0.3em}
{\footnotesize \emph{Notes:} This table reports standardized effect sizes (SDE) to facilitate cross-study comparison of treatment effect magnitudes. For continuous treatments, SDE $= \hat{\beta} \times \text{SD}(X) / \text{SD}(Y)$, which gives the effect of a one-standard-deviation change in the treatment variable, measured in standard deviations of the outcome. SD($Y$) and SD($X$) are unconditional standard deviations from the full analysis sample.

\textbf{Research question:} Does capping the SALT deduction at \$10,000 (TCJA 2018) and subsequently raising the cap to \$40,000 (OBBB 2025) affect local house prices? \textbf{Treatment:} Continuous; SALT Bite = max(0, avg SALT per itemizer $-$ \$10,000) / \$10,000, measured in \$10K units. \textbf{Data:} Zillow ZHVI (zip-code monthly, 2012--2026) merged with 2017 IRS SOI zip-code SALT data. 3,912,099 zip-month observations. \textbf{Method:} Continuous-treatment difference-in-differences with zip and metro$\times$month fixed effects, state-clustered standard errors. \textbf{Sample:} 25,303 zip codes observed monthly over 2012--2026.

Classification labels refer to the magnitude of the standardized point estimate, not to statistical significance. ``Null'' denotes a near-zero effect size ($|\text{SDE}| < 0.005$), not a failure to reject a null hypothesis.}
\end{table}
